\chapter{The Final Round and Trading Games}\label{ch:iv:the-final-round-and-trading-games}

Six candidates sit around a table with chips and a deck of cards; the interviewer asks for a market on the sum
of the next three cards. One quotes 17 at 23 and never moves it; one tightens after every trade and is picked
off twice by the same player; one says ``I am lifting the offer because two high cards are gone.'' The
assessor behind them writes down three different things, none of which is the chip count. A \omterm{def:iv:the-final-round-and-trading-games:final}{final round} is
the most expensive stage of the funnel, and each of its parts records something specific.

\section{The final round}

\begin{definition}[Final round, superday]\label{def:iv:the-final-round-and-trading-games:final}
The \emph{final round}\index{final round} is the last stage of an \omterm{def:iv:what-each-interview-tests-by-role:loop}{interview loop} before the decision: several
interviews on the same day or over a few days, each measuring different abilities, often with a game, a \omterm{def:iv:the-final-round-and-trading-games:group}{group
exercise} or a presentation. A \emph{superday}\index{superday} is the name some banks give to a final round
held on a single day with several interviewers (\cref{dat:iv:the-process-by-firm-type:published}).
\end{definition}

The \omterm{def:iv:the-final-round-and-trading-games:final}{final round} is designed so that each interview measures a different part of the role's analysis
(\cref{ch:iv:what-each-interview-tests-by-role}), and its decision is taken by a committee or a lead from the
written scores. Two consequences follow for the candidate. The day is long, and the fourth interview counts as
much as the first: energy is a resource to plan. And no single interview has to be perfect; a weak interview
among strong ones is usually recoverable, whereas a pattern (no checks, no numbers, no questions) is not.

\begin{method}[Planning a final-round day]\label{met:iv:the-final-round-and-trading-games:day}
\begin{enumerate}
\item Ask in advance for the schedule: how many interviews, of what kinds, with whom, and whether there is a
game or a presentation.
\item Prepare one-minute answers for the three questions every interviewer may ask: your background, why this
role, and one piece of work you know deeply.
\item Treat each interview as a fresh start; do not carry a bad one into the next room.
\item Eat and drink between interviews, even when told there is no time.
\item Write down, after each interview, the questions asked: the debrief with the recruiter goes better with
them, and so does the next process.
\end{enumerate}
\end{method}

\section{Mock trading and betting games}

\begin{definition}[Trading game]\label{def:iv:the-final-round-and-trading-games:game}
A \emph{trading game}\index{trading game} is an interview exercise in which candidates make markets on, bet
on or trade an uncertain quantity whose distribution can be reasoned about (dice, cards, an estimation
question), against the interviewer or each other, while information is revealed.
\end{definition}

The mechanics of making a market, its width and centre, and updating on trades, are in One Quant Book~2,
chapter~30, and the betting and sizing arithmetic in its chapter~29; this book's
\cref{ch:iv:betting-and-market-making-games} is the question bank for them. What matters in a \omterm{def:iv:the-final-round-and-trading-games:final}{final round} is
what the assessor records, which is behaviour: whether the candidate's price is near fair value and moves when
information arrives, whether the width reflects the uncertainty, whether the candidate notices being traded
against by someone who knows more, and whether positions and risk are tracked.

\begin{example}[The three-card market]\label{ex:iv:the-final-round-and-trading-games:cards}
Cards count 1 (ace) to 13 (king), four of each, 52 in all. The sum of three cards drawn without replacement
has mean $3 \times 7 = 21$ and variance $3 \times 14 \times \tfrac{49}{51} \approx 40.4$ (each card has
variance $(13^2 - 1)/12 = 14$, reduced by the finite-population factor), so a standard deviation of about 6.35;
an exact enumeration puts 90\% of the outcomes between 10 and 32. A market of 19 at 23 is centred and about a
third of a standard deviation wide on each side. When a card is shown, the fair value becomes that card plus
twice the mean of the 51 remaining cards: a king moves it to $13 + 2 \times 351/51 \approx 26.76$.
\Cref{fig:iv:the-final-round-and-trading-games:game} follows one game's fair value as its cards appear.
\end{example}

\begin{omfigure}
\begin{tikzpicture}
\begin{axis}[width=9.5cm,height=5.2cm,xlabel={\small cards revealed},ylabel={\small sum of the three cards},
  xmin=-0.4,xmax=3.4,ymin=5,ymax=35,xtick={0,1,2,3},
  tick label style={font=\footnotesize},grid=major,grid style={gray!20},table/col sep=comma]
\addplot[omDef,thick,mark=*,mark size=1.6pt,error bars/.cd,y dir=both,y explicit]
  table[x=revealed,y=fair,y error plus expr=\thisrow{hi}-\thisrow{fair},y error minus expr=\thisrow{fair}-\thisrow{lo}]
  {figdata/interviews/06-the-final-round-and-trading-games/game.csv};
\node[font=\scriptsize,anchor=south west] at (axis cs:1.05,28.5) {6 shown};
\node[font=\scriptsize,anchor=north west] at (axis cs:2.05,11) {then 3};
\node[font=\scriptsize,anchor=south west] at (axis cs:2.7,19) {then 7};
\end{axis}
\end{tikzpicture}

\omcaption{One three-card game: the fair value of the sum (dots) and the range holding 90\% of the outcomes
(bars) after 0, 1, 2 and 3 cards are shown. The cards were 6, 3 and 7 (a seeded draw). Exact enumeration;
data: \texttt{fig\_iv\_game.py}.}\label{fig:iv:the-final-round-and-trading-games:game}
\end{omfigure}

\begin{remark}[What the three candidates of the hook showed]\label{rem:iv:the-final-round-and-trading-games:hook}
A fixed market of 17 at 23 is centred at 20, below the fair 21, and never moves: it shows no updating. A
market tightened after every trade shows that the candidate treats trades as noise, whereas a player who
keeps trading on the same side is telling the market something (the adverse selection of One Quant Book~1,
chapter~1). Lifting an offer because two high cards are gone is conditional reasoning spoken aloud, which is
what the assessor came to see.
\end{remark}

\section{Group exercises and presentations}

\begin{definition}[Group exercise]\label{def:iv:the-final-round-and-trading-games:group}
A \emph{group exercise}\index{group exercise} is a final-round task given to several candidates at once (an
estimation, a case, a trading simulation with teams) and observed by assessors who score each candidate's
contribution, not the group's result.
\end{definition}

\omterm{def:iv:the-final-round-and-trading-games:group}{Group exercises} come from the assessment-centre tradition of personnel selection, in which several exercises
are scored on a few dimensions by trained assessors. A meta-analysis of 34 studies reduced the dimensions to
six (consideration of others, communication, drive, influencing others, organising and planning, problem
solving) and found criterion-related validities of 0.25 to 0.39 for them (Arthur, Day, McNelly and Edens,
2003). The candidate's lesson is that the group's answer is not scored; the candidate's part in reaching it
is. Talking most is not the same as contributing: structuring the problem, bringing a number, asking the quiet
member for theirs, and summarising the decision all score.

A presentation, often of the candidate's own past work or of a case prepared in the day, is scored on
structure, clarity and the response to challenge. The challenge is the point: an assessor who says ``your
sample is too small'' wants to see whether the candidate can quantify the concern, concede what is true and
defend what is not.

\section{What assessors write down}

The notes an assessor writes are short and concrete: the answer given, the time taken, whether a check was
made, how the candidate responded to new information, and one or two quotations. They are written against the
rubric's dimensions, and the debrief afterwards compares them across assessors before the decision. A
candidate can ask for feedback after the decision; many firms give a little, and what they give is usually a
paraphrase of those notes.

\section{Worked answers}

\begin{example}[A second game: how many suits]\label{ex:iv:the-final-round-and-trading-games:suits}
\emph{``Four cards will be dealt from a full deck. Make me a market on the number of different suits among them.''}
The value is 1, 2, 3 or 4. By indicators (\cref{ch:iv:probability-ii}), each suit is absent with probability
$\binom{39}{4}/\binom{52}{4} \approx 0.304$, so the expected number of suits is $4 \times (1 - 0.304) \approx 2.78$; all
four suits appear with probability $13^4/\binom{52}{4} \approx 0.11$. A market of 2.5 at 3.0 is centred. The first card
is shown and is a heart: by symmetry nothing changes, since some suit had to come first, and the expectation stays at
2.78 (computed from the three remaining cards: $1 + 3(1 - \binom{38}{3}/\binom{51}{3})$). The second card is also a
heart: now $1 + 3(1 - \binom{37}{2}/\binom{50}{2}) \approx 2.37$, and the market moves down to about 2.1 at 2.6. The
assessor scores the candidate who says aloud why the first card changed nothing and the second changed a lot.
\end{example}

\begin{example}[Answering ``your sample is too small'']\label{ex:iv:the-final-round-and-trading-games:challenge}
A candidate presents a strategy with an annualised Sharpe ratio of 1.5 over eighteen months of daily returns, and the
assessor says the sample is too small. The scoring response quantifies the objection before answering it: ``Measured on $T$ years, a
Sharpe ratio is uncertain by roughly $\sqrt{(1 + \mathrm{SR}^2/2)/T}$, here $\sqrt{2.125/1.5} \approx 1.19$. So
1.5 is only about 1.3 standard errors from zero: you are right that eighteen months cannot establish it. What I can
defend is the mechanism, which I tested on the other markets in the appendix, and what I'd want is a live test sized to
find out.'' Conceding the true part with a number and defending the rest is what the assessor writes down
(\cref{ch:iv:statistics}).
\end{example}

\section{Question bank}

\begin{interviewq}[$\star$ \iqroles{trader} \iqfirm{market maker}]\label{iq:iv:the-final-round-and-trading-games:1}
Make me a market on the sum of three cards drawn from a full deck, counting ace as 1 and king as 13. Give
your centre and your width, and say why.
\end{interviewq}

\begin{interviewq}[$\star$ \iqroles{trader, researcher, developer} \iqfirm{bank}]\label{iq:iv:the-final-round-and-trading-games:2}
You have a \omterm{def:iv:the-final-round-and-trading-games:final}{superday} of five interviews tomorrow, from 9:00 to 15:00. What do you do tonight, and how do you
manage the day itself?
\end{interviewq}

\begin{interviewq}[$\star$ \iqroles{trader} \iqfirm{market maker}]\label{iq:iv:the-final-round-and-trading-games:3}
In the three-card market, the first card is turned over and it is a king. Where is fair value now?
\end{interviewq}

\begin{interviewq}[$\star\star$ \iqroles{trader} \iqfirm{market maker}]\label{iq:iv:the-final-round-and-trading-games:4}
You quote 19 at 23 on the three-card sum. The interviewer has been shown one of the three cards and buys at
23. If the interviewer buys whenever that card is above 7, what is the fair value given the trade, and how do
you move your market?
\end{interviewq}

\begin{interviewq}[$\star\star$ \iqroles{trader, risk} \iqfirm{proprietary firm}]\label{iq:iv:the-final-round-and-trading-games:5}
You start a betting game with 100 chips. Ten times, you may bet any part of your chips at even money on a coin
that lands heads with probability 0.6. What fraction do you bet each time if you want to maximise the
expected logarithm of your final chips? What is the expected growth per bet, and what does it imply for the
typical final stack?
\end{interviewq}

\begin{interviewq}[$\star\star$ \iqroles{trader, researcher} \iqfirm{market maker}]\label{iq:iv:the-final-round-and-trading-games:6}
In a \omterm{def:iv:the-final-round-and-trading-games:group}{group exercise}, five candidates must estimate the number of trading days on which a large equity index
moved more than 2\% over the last twenty years, and agree one answer in fifteen minutes. What role do you take,
and what do you actually say in the first two minutes?
\end{interviewq}

\begin{interviewq}[$\star\star$ \iqroles{trader} \iqfirm{market maker}]\label{iq:iv:the-final-round-and-trading-games:7}
You are long three contracts on the three-card sum at an average price of 24. Two cards are shown: 5 and 6.
What is your expected P\&L? Someone bids 18. Do you sell?
\end{interviewq}

\begin{interviewq}[$\star\star\star$ \iqroles{trader, risk} \iqfirm{market maker}]\label{iq:iv:the-final-round-and-trading-games:8}
You make a market centred at 21 on the three-card sum. A player who has seen one of the three cards trades
whenever your bid or offer is wrong in their favour. What is your expected loss per trade opportunity with a
half-width of 0? Of 2? What half-width stops this player from ever trading?
\end{interviewq}

\begin{interviewq}[$\star\star\star$ \iqroles{trader, researcher} \iqfirm{proprietary firm}]\label{iq:iv:the-final-round-and-trading-games:9}
A deck of 30 red and 22 black cards is shuffled and turned over one card at a time. At any moment before the
last card you may say ``next'', and you win if the next card is red; if you never say it, the last card
decides. What stopping rule maximises your chance of winning, and what is the chance?
\end{interviewq}

\begin{interviewq}[$\star\star\star$ \iqroles{trader, researcher} \iqfirm{market maker}]\label{iq:iv:the-final-round-and-trading-games:10}
You are the assessor for the hook's three candidates. Write the one-line note you would record for each, and
say which of them you would pass and why.
\end{interviewq}

\begin{omsources}
\item W.~Arthur, E.~A.~Day, T.~L.~McNelly and P.~S.~Edens, ``A meta-analysis of the criterion-related validity
of assessment center dimensions'', \emph{Personnel Psychology} 56(1), 2003, 125--153.
\item International Task Force on Assessment Center Guidelines, ``Guidelines and ethical considerations for
assessment center operations'', \emph{International Journal of Selection and Assessment} 17(3), 2009,
243--253.
\item J.~L.~Kelly, ``A new interpretation of information rate'', \emph{Bell System Technical Journal} 35(4),
1956, 917--926.
\item One Quant Book~2, chapters 29--30 (expected value, Kelly, making markets); One Quant Book~4, chapter~1
(martingales and stopping).
\end{omsources}
